\documentclass[11pt]{article}
\usepackage{amsmath, amssymb}
\usepackage[round,authoryear]{natbib}
\usepackage{hyperref}
\usepackage{xcolor}
\usepackage{booktabs}
\usepackage{graphicx}
\usepackage{caption}
\usepackage{float}
\usepackage{longtable}
\usepackage{array}
\usepackage{mathptmx}
\usepackage[T1]{fontenc}
\usepackage[margin=1in]{geometry}
\usepackage{setspace}
\setstretch{1.15}
\usepackage{threeparttable}
\captionsetup{labelfont=bf, textfont=it}

\hypersetup{
    colorlinks=true,
    linkcolor=blue,
    citecolor=blue,
    urlcolor=blue
}

\title{Online Appendix: Simulation Results\\[0.5em]
\large An Adaptive-Hybrid New-Keynesian Model}
\author{Bruno Cittolin Smaniotto}
\date{}

\begin{document}

\maketitle

\begin{abstract}
This online appendix documents the numerical results underlying the claims in ``An Adaptive-Hybrid New-Keynesian Model.'' Each table presents simulation outputs with corresponding parameter configurations and references to the generating code. All results can be replicated using the code available in the replication package.
\end{abstract}

\tableofcontents
\newpage

%%%%%%%%%%%%%%%%%%%%%%%%%%%%%%%%%%%%%%%%%%%%%%%%%%%%%%%%%%%%%%%%%%%%%%%%%%%%%%%
\section{Introduction}

This document serves as a comprehensive reference for the numerical claims made in the main text. For each simulation result, we report:
\begin{itemize}
    \item The specific numerical values cited in the manuscript
    \item The parameter configuration used to generate them
    \item The code file that produces the result
\end{itemize}

All simulations use the baseline calibration from Table 1 of the main text unless otherwise noted. The complete replication package, including all code and data, is available at \url{https://github.com/brunosmaniotto/An-Adaptive-Hybrid-New-Keynesian-Model}.

%%%%%%%%%%%%%%%%%%%%%%%%%%%%%%%%%%%%%%%%%%%%%%%%%%%%%%%%%%%%%%%%%%%%%%%%%%%%%%%
\section{Baseline Calibration}
\label{sec:calibration}

Table~\ref{tab:baseline_params} presents the baseline parameter values used throughout the simulations. These are defined in \texttt{code/models/parameters.py}.

\begin{table}[htbp]
\centering
\caption{Baseline Parameter Values}
\label{tab:baseline_params}
\begin{threeparttable}
\begin{tabular}{llcl}
\toprule
Parameter & Description & Value & Source \\
\midrule
\multicolumn{4}{l}{\textit{Structural Parameters}} \\
$\beta$ & Discount factor & 0.99 & Gal\'i (2015) \\
$\sigma$ & Intertemporal elasticity of substitution & 1.0 & Log utility \\
$\kappa$ & Phillips curve slope & 0.024 & Gal\'i \& Gertler (1999) \\
$\phi_\pi$ & Taylor rule: inflation response & 1.5 & Taylor (1993) \\
$\phi_y$ & Taylor rule: output response & 0.125 & Taylor (1993) \\
$\pi^*$ & Inflation target (quarterly) & 0.005 & 2\% annual \\
$\rho_u$ & Cost-push shock persistence & 0.5 & -- \\
$\rho_r$ & Natural rate persistence & 0.9 & -- \\
\midrule
\multicolumn{4}{l}{\textit{Behavioral Parameters}} \\
$\lambda$ & Fraction of FIRE agents & 0.35 & Coibion \& Gorodnichenko (2015) \\
$\eta$ & Learning speed & 0.10 & Malmendier \& Nagel (2016) \\
$k$ & Memory window (quarters) & 3 & -- \\
$\varepsilon$ & Simplicity bias threshold & $10^{-4}$ & Appendix B \\
\bottomrule
\end{tabular}
\begin{tablenotes}
\small
\item \textit{Code reference:} \texttt{code/models/parameters.py}
\end{tablenotes}
\end{threeparttable}
\end{table}

%%%%%%%%%%%%%%%%%%%%%%%%%%%%%%%%%%%%%%%%%%%%%%%%%%%%%%%%%%%%%%%%%%%%%%%%%%%%%%%
\section{Section 2: The Model}
\label{sec:model_results}

\subsection{Figure 1: Learning Mechanism}

The learning mechanism illustration uses a stylized simulation with the following configuration:

\begin{table}[htbp]
\centering
\caption{Figure 1 Simulation Parameters and Results}
\label{tab:fig1_results}
\begin{threeparttable}
\begin{tabular}{lcc}
\toprule
Quantity & Value & Manuscript Reference \\
\midrule
Learning rate $\eta$ & 0.10 & Figure~1 caption \\
Initial credibility $\theta_0$ & 1.0 & -- \\
Credibility during stable periods & $\approx 0.95$ & Section~2.1, Figure~1 discussion \\
Credibility after persistent shock (Q35) & $\approx 0.40$ & Section~2.1, Figure~1 discussion \\
Credibility after 25Q recovery & $\approx 0.70$ & Section~2.1, Figure~1 discussion \\
\bottomrule
\end{tabular}
\begin{tablenotes}
\small
\item \textit{Code reference:} \texttt{code/simulations/section\_2/learning\_mechanism.py}, \texttt{code/plotting/section\_2/figure\_01\_learning\_mechanism.py}
\item \textit{Output:} \texttt{output/figures/section\_2/figure\_01\_learning\_mechanism.png}
\end{tablenotes}
\end{threeparttable}
\end{table}

\subsection{Figure 2: Roles of $k$ and $\varepsilon$}

Panel (a) demonstrates the re-anchoring problem when $\varepsilon = 0$. Panel (b) compares memory window lengths.

\begin{table}[htbp]
\centering
\caption{Figure 2 Results: Memory Window Comparison}
\label{tab:fig2_results}
\begin{threeparttable}
\begin{tabular}{lccc}
\toprule
Configuration & De-anchoring Trigger & Recovery Possible & Notes \\
\midrule
\multicolumn{4}{l}{\textit{Panel (a): Simplicity Bias}} \\
$\varepsilon = 0$ & Single-quarter spike & No & Credibility stuck at 0 \\
$\varepsilon = 10^{-4}$ & Persistent deviation & Yes & Gradual recovery \\
\midrule
\multicolumn{4}{l}{\textit{Panel (b): Memory Window}} \\
$k = 1$ & 2-quarter spike & Immediate & High volatility \\
$k = 3$ & 12-quarter deviation & Delayed & More stable \\
\bottomrule
\end{tabular}
\begin{tablenotes}
\small
\item \textit{Code reference:} \texttt{code/simulations/section\_2/k\_epsilon\_roles.py}, \texttt{code/plotting/section\_2/figure\_02\_k\_epsilon\_roles.py}
\item \textit{Output:} \texttt{output/figures/section\_2/figure\_02\_k\_epsilon\_roles.png}
\end{tablenotes}
\end{threeparttable}
\end{table}

\subsection{Figure 3: Effect of Rational Agent Fraction}

Simulations compare $\lambda \in \{0, 0.3, 1.0\}$ under single and repeated shocks.

\begin{table}[htbp]
\centering
\caption{Figure 3 Simulation Configuration}
\label{tab:fig3_config}
\begin{threeparttable}
\begin{tabular}{lcc}
\toprule
Parameter & Single Shock & Repeated Shocks \\
\midrule
Shock magnitude & 0.5\% & 0.75\% \\
Shock persistence $\rho$ & 0.8 & -- \\
Shock duration & -- & 12 quarters \\
$\lambda$ values & \multicolumn{2}{c}{$\{0, 0.3, 1.0\}$} \\
\bottomrule
\end{tabular}
\begin{tablenotes}
\small
\item \textit{Key finding:} Mixed population ($\lambda = 0.3$) produces highest peak inflation.
\item \textit{Code reference:} \texttt{code/simulations/section\_2/lambda\_comparison.py}, \texttt{code/plotting/section\_2/figure\_03\_lambda\_comparison.py}
\item \textit{Output:} \texttt{output/figures/section\_2/figure\_03\_lambda\_comparison.png}
\end{tablenotes}
\end{threeparttable}
\end{table}

\subsection{Table 2: Transitory vs.\ Persistent Shocks}

This table demonstrates that shock persistence, not magnitude, triggers de-anchoring.

\begin{table}[htbp]
\centering
\caption{Transitory vs.\ Persistent Shocks: Simulation Results}
\label{tab:transitory_persistent}
\begin{threeparttable}
\begin{tabular}{ccccc}
\toprule
Persistence ($\rho$) & Peak $\pi$ (FIRE) & Peak $\pi$ (Adaptive) & Difference & Min $\theta$ \\
\midrule
0.0 & 4.33\% & 4.36\% & 0.03\% & 1.00 \\
0.4 & 7.00\% & 4.27\% & $-$2.73\% & 1.00 \\
0.85 & 22.78\% & 3.76\% & $-$19.02\% & 0.59 \\
\bottomrule
\end{tabular}
\begin{tablenotes}
\small
\item \textit{Configuration:} 4.5\% annualized shock, baseline parameters ($\eta = 0.10$, $\phi_\pi = 1.5$, $k = 3$, $\varepsilon = 10^{-4}$)
\item \textit{Code reference:} \texttt{code/tables/section\_2/table\_02\_transitory\_persistent.py}
\end{tablenotes}
\end{threeparttable}
\end{table}

\subsection{Table 3: The Paradox of Faster Learning}

Faster learning ($\eta$) leads to worse inflation outcomes under persistent shocks.

\begin{table}[htbp]
\centering
\caption{The Paradox of Faster Learning: Simulation Results}
\label{tab:paradox_learning}
\begin{threeparttable}
\begin{tabular}{ccccc}
\toprule
Learning Speed ($\eta$) & Peak $\Delta\pi$ & Min $\theta$ & Quarters to Target & Final $\theta$ \\
\midrule
0.10 & 6.68\% & 0.122 & 28 & 1.000 \\
0.20 & 8.35\% & 0.000 & 96 & 0.995 \\
0.30 & 14.75\% & 0.000 & $>$115 & 0.000 \\
\bottomrule
\end{tabular}
\begin{tablenotes}
\small
\item \textit{Configuration:} 8\% annualized persistent shock ($\rho = 0.85$), baseline parameters ($\phi_\pi = 1.5$, $k = 3$, $\varepsilon = 10^{-4}$)
\item \textit{Code reference:} \texttt{code/tables/section\_2/table\_03\_paradox\_learning.py}
\end{tablenotes}
\end{threeparttable}
\end{table}

\subsection{Figure 4: Shock Regimes}

Three distinct shock environments are simulated to illustrate regime-dependent dynamics.

\begin{table}[htbp]
\centering
\caption{Figure 4: Shock Regime Configurations and Results}
\label{tab:shock_regimes}
\begin{threeparttable}
\begin{tabular}{lccc}
\toprule
Regime & Shock $\sigma_u$ & Key Result & $\theta$ Dynamics \\
\midrule
Great Moderation & 0.0006 & Inflation stable near 2\% & $\theta$ fluctuates around 0.5 \\
1970s-Style & Large, volatile & Peaks near 10\% & $\theta$ erodes from 0.9 \\
Credibility Cycle & Variable & Full cycle & Recovery $\approx$ 15 years \\
\bottomrule
\end{tabular}
\begin{tablenotes}
\small
\item \textit{Code reference:} \texttt{code/simulations/section\_2/shock\_regimes.py}, \texttt{code/plotting/section\_2/figure\_04\_shock\_regimes.py}
\item \textit{Output:} \texttt{output/figures/section\_2/figure\_04\_shock\_regimes.png}
\end{tablenotes}
\end{threeparttable}
\end{table}

%%%%%%%%%%%%%%%%%%%%%%%%%%%%%%%%%%%%%%%%%%%%%%%%%%%%%%%%%%%%%%%%%%%%%%%%%%%%%%%
\section{Section 2: Empirical Evidence on Persistence}
\label{sec:persistence}

\subsection{Figure 5: Time-Varying Inflation Persistence}

Rolling 40-quarter AR(1) regressions on U.S.\ CPI inflation (1960Q1--2024Q4).

\begin{table}[htbp]
\centering
\caption{Inflation Persistence by Monetary Policy Regime (corresponds to Table 4 in main text)}
\label{tab:persistence_regimes}
\begin{threeparttable}
\begin{tabular}{lcccccc}
\toprule
Regime & Period & Mean $\hat{\rho}$ & Std.\ Dev. & $N$ & $F$-stat & $p$-value \\
\midrule
Great Inflation & 1968--1979 & 0.729 & 0.042 & 22 & 13.13 & $<$0.001 \\
Volcker Era & 1980--1984 & 0.671 & 0.080 & 20 & 23.77 & $<$0.001 \\
Great Moderation & 1985--2007 & 0.255 & 0.319 & 92 & 33.55 & $<$0.001 \\
Post-Crisis & 2008--2019 & $-$0.045 & 0.105 & 48 & 1.93 & 0.147 \\
COVID Era & 2020--2024 & 0.392 & 0.221 & 20 & -- & -- \\
\bottomrule
\end{tabular}
\begin{tablenotes}
\small
\item \textit{Method:} Rolling AR(1) with 40-quarter window; Chow tests for structural breaks
\item \textit{Data:} FRED CPI series (CPIAUCSL)
\item \textit{Code reference:} \texttt{code/empirical/section\_2/persistence\_analysis.py}, \texttt{code/plotting/section\_2/figure\_05\_persistence.py}
\item \textit{Output:} \texttt{output/figures/section\_2/figure\_05\_persistence.png}, \texttt{output/empirical/section\_2/persistence\_estimates.csv}
\end{tablenotes}
\end{threeparttable}
\end{table}

%%%%%%%%%%%%%%%%%%%%%%%%%%%%%%%%%%%%%%%%%%%%%%%%%%%%%%%%%%%%%%%%%%%%%%%%%%%%%%%
\section{Section 3: Applications and Policy Implications}
\label{sec:applications}

\subsection{Figure 6: Oil Shocks and Initial Credibility}

Identical shocks produce dramatically different outcomes depending on initial credibility.

\begin{table}[htbp]
\centering
\caption{Figure 6: Oil Shocks Simulation Results}
\label{tab:oil_shocks}
\begin{threeparttable}
\begin{tabular}{lcc}
\toprule
Quantity & Low $\theta_0$ (0.2) & High $\theta_0$ (1.0) \\
\midrule
Shock magnitude & \multicolumn{2}{c}{100 bp cost-push} \\
Policy rule & \multicolumn{2}{c}{$\phi_\pi = 1.5$} \\
\midrule
Peak inflation & 8.2\% & 5.5\% \\
Peak difference & \multicolumn{2}{c}{2.7 pp} \\
Quarters above 3\% & 12 & 4 \\
Persistence ratio & \multicolumn{2}{c}{$3.0\times$} \\
Final credibility & 0.11 & 1.00 \\
\bottomrule
\end{tabular}
\begin{tablenotes}
\small
\item \textit{Code reference:} \texttt{code/simulations/section\_3/oil\_shocks.py}, \texttt{code/plotting/section\_3/figure\_06\_oil\_shocks.py}
\item \textit{Output:} \texttt{output/figures/section\_3/figure\_06\_oil\_shocks.png}
\end{tablenotes}
\end{threeparttable}
\end{table}

\subsection{Figure 7: The Great Moderation}

High inherited credibility makes the adaptive model behave like standard NK.

\begin{table}[htbp]
\centering
\caption{Figure 7: Great Moderation Simulation Results}
\label{tab:great_moderation}
\begin{threeparttable}
\begin{tabular}{lc}
\toprule
Quantity & Value \\
\midrule
Initial credibility $\theta_0$ & 0.95 \\
Shock volatility $\sigma$ & 0.0006 \\
Shock persistence $\rho$ & 0.3 \\
\midrule
Correlation (Adaptive vs.\ NK) & 0.94 \\
Maximum difference & $<$ 0.2 pp \\
Credibility throughout & $\approx 1.0$ \\
\bottomrule
\end{tabular}
\begin{tablenotes}
\small
\item \textit{Code reference:} \texttt{code/simulations/section\_3/great\_moderation.py}, \texttt{code/plotting/section\_3/figure\_07\_great\_moderation.py}
\item \textit{Output:} \texttt{output/figures/section\_3/figure\_07\_great\_moderation.png}
\end{tablenotes}
\end{threeparttable}
\end{table}

\subsection{Figure 8: The Missing Disinflation}

Anchored expectations absorb the demand collapse, preventing deflation.

\begin{table}[htbp]
\centering
\caption{Figure 8: Missing Disinflation Simulation Results}
\label{tab:missing_disinflation}
\begin{threeparttable}
\begin{tabular}{lcc}
\toprule
Quantity & Standard NK & Adaptive Model \\
\midrule
Initial credibility $\theta_0$ & -- & 0.95 \\
Shock magnitude & \multicolumn{2}{c}{$-$50 bp} \\
Shock persistence $\rho$ & \multicolumn{2}{c}{0.7} \\
\midrule
Minimum inflation & $-$3.7\% & $-$0.6\% \\
Difference & \multicolumn{2}{c}{3.1 pp} \\
\bottomrule
\end{tabular}
\begin{tablenotes}
\small
\item \textit{Code reference:} \texttt{code/simulations/section\_3/missing\_disinflation.py}, \texttt{code/plotting/section\_3/figure\_08\_missing\_disinflation.py}
\item \textit{Output:} \texttt{output/figures/section\_3/figure\_08\_missing\_disinflation.png}
\end{tablenotes}
\end{threeparttable}
\end{table}

\subsection{Figure 9: Post-Pandemic Inflation}

Two-phase dynamics, with initial dampening followed by persistence as credibility erodes.

\begin{table}[htbp]
\centering
\caption{Figure 9: Post-Pandemic Simulation Results}
\label{tab:post_pandemic}
\begin{threeparttable}
\begin{tabular}{lcc}
\toprule
Quantity & Standard NK & Adaptive Model \\
\midrule
Initial credibility $\theta_0$ & -- & 0.6 \\
Policy response $\phi_\pi$ & \multicolumn{2}{c}{1.1 (weak)} \\
\midrule
Peak inflation & 16.7\% & 7.8\% \\
Return to target & 4.5 years & 8.0 years \\
Crossover point & \multicolumn{2}{c}{3.5 years} \\
Final credibility & -- & 0.10 \\
\bottomrule
\end{tabular}
\begin{tablenotes}
\small
\item \textit{Note:} Crossover marks when adaptive model costs exceed initial dampening benefits.
\item \textit{Code reference:} \texttt{code/simulations/section\_3/post\_pandemic.py}, \texttt{code/plotting/section\_3/figure\_09\_post\_pandemic.py}
\item \textit{Output:} \texttt{output/figures/section\_3/figure\_09\_post\_pandemic.png}
\end{tablenotes}
\end{threeparttable}
\end{table}

\subsection{Figure 10: Long and Variable Lags}

Transmission speed varies systematically with credibility state.

\begin{table}[htbp]
\centering
\caption{Figure 10: Transmission Lags by Credibility}
\label{tab:transmission_lags}
\begin{threeparttable}
\begin{tabular}{lccc}
\toprule
Initial $\theta_0$ & Peak Inflation & Quarters to Target & Relative Speed \\
\midrule
0.2 (Low) & 8.7\% & 14 & Baseline \\
0.5 (Medium) & 7.5\% & 12 & 17\% faster \\
0.8 (High) & 6.3\% & 10 & 40\% faster \\
\midrule
Shock & \multicolumn{3}{c}{0.75 pp cost-push, $\rho = 0.75$} \\
\bottomrule
\end{tabular}
\begin{tablenotes}
\small
\item \textit{Code reference:} \texttt{code/simulations/section\_3/transmission\_lags.py}, \texttt{code/plotting/section\_3/figure\_10\_transmission\_lags.py}
\item \textit{Output:} \texttt{output/figures/section\_3/figure\_10\_transmission\_lags.png}
\end{tablenotes}
\end{threeparttable}
\end{table}

\subsection{Figure 11: Phillips Curve Estimation}

Simulated data reveal apparent parameter instability arising from credibility variation.

\begin{table}[htbp]
\centering
\caption{Figure 11: Phillips Curve Estimates by Credibility Regime}
\label{tab:pc_estimation}
\begin{threeparttable}
\begin{tabular}{lccc}
\toprule
Credibility Regime & AR(1) $\hat{\rho}$ & Standard $\hat{\kappa}$ & Hybrid $\hat{\kappa}$ \\
\midrule
Low ($\theta = 0.2$) & 0.97 & 0.062 & 0.016 \\
Medium ($\theta = 0.5$) & 0.91 & 0.038 & 0.016 \\
High ($\theta = 0.95$) & 0.79 & 0.026 & 0.016 \\
\bottomrule
\end{tabular}
\begin{tablenotes}
\small
\item \textit{Configuration:} 200 quarters, shocks $\sigma = 0.002$, $\rho = 0.3$
\item \textit{Note:} Hybrid specification includes $\pi_{t-1}$; standard specification omits it.
\item \textit{Code reference:} \texttt{code/simulations/section\_3/phillips\_curve.py}, \texttt{code/plotting/section\_3/figure\_11\_phillips\_curve.py}
\item \textit{Output:} \texttt{output/figures/section\_3/figure\_11\_phillips\_curve.png}
\end{tablenotes}
\end{threeparttable}
\end{table}

\subsection{Figure 12: State-Dependent Optimal Policy}

Optimal policy aggressiveness varies inversely with credibility.

\begin{table}[htbp]
\centering
\caption{Figure 12: Optimal Policy Coefficient by Initial Credibility}
\label{tab:optimal_policy}
\begin{threeparttable}
\begin{tabular}{lcccc}
\toprule
Initial $\theta_0$ & Optimal $\phi_\pi^*$ & Relative Loss & Taylor Principle \\
\midrule
0.05 & $\approx 3.0$ & 3.6$\times$ baseline & Required \\
0.35 & $\approx 2.0$ & 2.0$\times$ baseline & Required \\
0.75 & $\approx 1.0$ & 1.2$\times$ baseline & Not required \\
0.95 & $\approx 0.8$ & Baseline & Not required \\
\bottomrule
\end{tabular}
\begin{tablenotes}
\small
\item \textit{Method:} Grid search over $\phi_\pi \in [0.5, 5.0]$, 16-quarter simulations
\item \textit{Note:} Softmax smoothing ($\gamma = 25{,}000$) used for differentiable welfare function.
\item \textit{Code reference:} \texttt{code/simulations/section\_3/policy\_asymmetry.py}, \texttt{code/plotting/section\_3/figure\_12\_policy\_asymmetry.py}
\item \textit{Output:} \texttt{output/figures/section\_3/figure\_12\_policy\_asymmetry.png}
\end{tablenotes}
\end{threeparttable}
\end{table}

\subsection{Figure 13: The Credibility Buffer}

High credibility permits delayed policy response; low credibility does not.

\begin{table}[htbp]
\centering
\caption{Figure 13: Credibility Buffer Simulation Results}
\label{tab:credibility_buffer}
\begin{threeparttable}
\begin{tabular}{lcccc}
\toprule
 & \multicolumn{2}{c}{High $\theta_0$ (0.95)} & \multicolumn{2}{c}{Low $\theta_0$ (0.50)} \\
\cmidrule(lr){2-3} \cmidrule(lr){4-5}
 & Immediate & Delayed & Immediate & Delayed \\
\midrule
Peak inflation & Similar & Similar & Moderate & High \\
Welfare loss & $\approx$ Equal & $\approx$ Equal & Baseline & Elevated \\
De-anchoring & No & No & Mild & Severe \\
\bottomrule
\end{tabular}
\begin{tablenotes}
\small
\item \textit{Configuration:} Delay period = 8 quarters (interest rate held at neutral)
\item \textit{Code reference:} \texttt{code/simulations/section\_3/credibility\_buffer.py}, \texttt{code/plotting/section\_3/figure\_13\_credibility\_buffer.py}
\item \textit{Output:} \texttt{output/figures/section\_3/figure\_13\_credibility\_buffer.png}
\end{tablenotes}
\end{threeparttable}
\end{table}

%%%%%%%%%%%%%%%%%%%%%%%%%%%%%%%%%%%%%%%%%%%%%%%%%%%%%%%%%%%%%%%%%%%%%%%%%%%%%%%
\section{Section 4: Extensions}
\label{sec:extensions}

\subsection{Figure 14: Hyperinflation and the Three-Arm Model}

Uses baseline parameters ($\lambda = 0.35$, $\eta = 0.10$, $\kappa = 0.024$) with a single cost-push shock of $0.08$ at $t=5$ and persistence $\rho_u = 0.90$. The Taylor principle ($\phi_\pi > 1$ vs.\ $\phi_\pi \leq 1$) determines whether trend-following persists.

\begin{table}[htbp]
\centering
\caption{Figure 14: Three-Arm Model Results}
\label{tab:hyperinflation}
\begin{threeparttable}
\begin{tabular}{lcccc}
\toprule
Policy Stance & $\phi_\pi$ & Peak Inflation & Max TF Share & TF Persists \\
\midrule
Active & $1.5$ & 60.3\% & 0.19 & No \\
Boundary & $1.0$ & 61.6\% & 0.19 & Yes \\
Passive & $0.5$ & 62.9\% & 0.27 & Yes \\
\bottomrule
\end{tabular}
\begin{tablenotes}
\small
\item \textit{Key finding:} Under active policy ($\phi_\pi > 1$), trend-following activates transiently but returns to zero. Under passive policy ($\phi_\pi \leq 1$), trend-following persists and can generate explosive dynamics.
\item \textit{Code reference:} \texttt{code/simulations/section\_4/hyperinflation.py}, \texttt{code/plotting/section\_4/figure\_14\_hyperinflation.py}
\item \textit{Output:} \texttt{output/figures/section\_4/figure\_14\_hyperinflation.png}
\end{tablenotes}
\end{threeparttable}
\end{table}

\subsection{Figure 15: Brazil---Counterfactual vs.\ Real Plan}

The URV re-anchored expectations before currency reform.

\begin{table}[htbp]
\centering
\caption{Figure 15: Brazil Simulation Results}
\label{tab:brazil}
\begin{threeparttable}
\begin{tabular}{lcc}
\toprule
Quantity & Counterfactual & Real Plan \\
\midrule
Initial state & $\theta_{CB} = 0.05$, $\theta_{TF} = 0.80$ & Same + URV phase \\
Shock & \multicolumn{2}{c}{2\% for 4 quarters} \\
Policy & \multicolumn{2}{c}{$\phi_\pi = 1.5$} \\
\midrule
Peak inflation (annualized) & $>$2000\% & 18\% \\
URV re-anchoring & -- & $\theta_{CB}$: 5\% $\to$ 41\% in 4Q \\
Final credibility & $\theta_{CB} \to 0$ & $\theta_{CB} \to 1$ \\
Welfare improvement & -- & $>$90\% \\
\bottomrule
\end{tabular}
\begin{tablenotes}
\small
\item \textit{Code reference:} \texttt{code/simulations/section\_4/brazil.py}, \texttt{code/plotting/section\_4/figure\_15\_brazil.py}
\item \textit{Output:} \texttt{output/figures/section\_4/figure\_15\_brazil.png}
\end{tablenotes}
\end{threeparttable}
\end{table}

\subsection{Figure 16: Long Memory and Inflation Trauma}

Exponentially discounted memory ($\delta = 0.8$) captures lasting scars.

\begin{table}[htbp]
\centering
\caption{Figure 16: Long Memory Simulation Results (20-Year Histories)}
\label{tab:long_memory}
\begin{threeparttable}
\begin{tabular}{llcc}
\toprule
Scenario & Description & Final $\theta$ & Peak $\pi$ \\
\midrule
Stable & At 2\% target throughout & 1.00 & 8.2\% \\
Distant trauma & 5y hyperinflation, then 15y stable & 0.47 & 12.1\% \\
Chronic (current) & 10y stable, then 10y at 7.5\% & 0.02 & 23.3\% \\
\bottomrule
\end{tabular}
\begin{tablenotes}
\small
\item \textit{Key finding:} Even 15 years after hyperinflation, credibility only partially recovers ($\theta \approx 0.5$).
\item \textit{Code reference:} \texttt{code/simulations/section\_4/long\_memory.py}, \texttt{code/plotting/section\_4/figure\_16\_long\_memory.py}
\item \textit{Output:} \texttt{output/figures/section\_4/figure\_16\_long\_memory.png}
\end{tablenotes}
\end{threeparttable}
\end{table}

\subsection{Figure 17: Japan---Anchored Below Target}

Symmetric de-anchoring occurs when inflation persistently undershoots.

\begin{table}[htbp]
\centering
\caption{Figure 17: Japan Simulation Results}
\label{tab:japan}
\begin{threeparttable}
\begin{tabular}{lcc}
\toprule
Inflation Scenario & Deviation from 2\% Target & Final $\theta$ \\
\midrule
At target (2\%) & 0 pp & 1.00 \\
Close (1\%) & $-$1.0 pp & High (within $\varepsilon$) \\
Japan (0\%) & $-$2.0 pp & $\approx 0$ \\
\bottomrule
\end{tabular}
\begin{tablenotes}
\small
\item \textit{Key finding:} 2 pp persistent miss causes credibility collapse regardless of direction.
\item \textit{Code reference:} \texttt{code/simulations/section\_4/japan.py}, \texttt{code/plotting/section\_4/figure\_17\_japan.py}
\item \textit{Output:} \texttt{output/figures/section\_4/figure\_17\_japan.png}
\end{tablenotes}
\end{threeparttable}
\end{table}

%%%%%%%%%%%%%%%%%%%%%%%%%%%%%%%%%%%%%%%%%%%%%%%%%%%%%%%%%%%%%%%%%%%%%%%%%%%%%%%
\section{Appendix Results}
\label{sec:appendix_results}

\subsection{Figure B1: The Re-Anchoring Problem}

Without the simplicity bias ($\varepsilon = 0$), credibility cannot recover.

\begin{table}[htbp]
\centering
\caption{Figure B1: Re-Anchoring Problem Results}
\label{tab:reanchoring}
\begin{threeparttable}
\begin{tabular}{lcc}
\toprule
Configuration & $\varepsilon = 0$ & $\varepsilon = 10^{-4}$ \\
\midrule
Shock & \multicolumn{2}{c}{2\%, $\rho = 0.8$} \\
Learning rate & \multicolumn{2}{c}{$\eta = 0.10$} \\
\midrule
De-anchoring & Identical & Identical \\
Recovery after stabilization & No & Yes \\
Final $\theta$ (post-shock) & $\approx 0$ & Gradual recovery \\
\bottomrule
\end{tabular}
\begin{tablenotes}
\small
\item \textit{Code reference:} \texttt{code/simulations/appendix\_B/reanchoring.py}, \texttt{code/plotting/appendix\_B/figure\_B1\_reanchoring.py}
\item \textit{Output:} \texttt{output/figures/appendix\_B/figure\_B1\_reanchoring.png}
\end{tablenotes}
\end{threeparttable}
\end{table}

\subsection{Appendix D: Robustness---Kalman Filter Comparison}

State-space estimation validates rolling-window persistence estimates.

\begin{table}[htbp]
\centering
\caption{Table A1: Kalman TVP vs.\ Rolling Window Persistence}
\label{tab:kalman_comparison}
\begin{threeparttable}
\begin{tabular}{lccc}
\toprule
Period & Kalman $\bar{\rho}$ & Rolling $\bar{\rho}$ & Difference \\
\midrule
Great Inflation (1968--1979) & 0.703 & 0.729 & $-$0.026 \\
Volcker Era (1980--1984) & 0.490 & 0.671 & $-$0.181 \\
Great Moderation (1985--2007) & 0.145 & 0.255 & $-$0.109 \\
Post-Crisis (2008--2019) & 0.009 & $-$0.045 & $+$0.054 \\
COVID Era (2020--2024) & 0.352 & 0.392 & $-$0.040 \\
\midrule
Correlation & \multicolumn{3}{c}{0.86} \\
\bottomrule
\end{tabular}
\begin{tablenotes}
\small
\item \textit{Kalman estimates:} $\hat{\sigma}^2_\varepsilon = 5.12$, $\hat{\sigma}^2_\omega = 0.0036$, signal-to-noise = 0.0007
\item \textit{Code reference:} \texttt{code/empirical/appendix\_D/kalman\_comparison.py}
\item \textit{Output:} \texttt{output/tables/appendix\_D/table\_D1.csv}
\end{tablenotes}
\end{threeparttable}
\end{table}

\subsection{Appendix D: Window Size Sensitivity}

Persistence estimates are stable across window sizes.

\begin{table}[htbp]
\centering
\caption{Table A2: Rolling Persistence by Window Size}
\label{tab:window_sensitivity}
\begin{threeparttable}
\begin{tabular}{lcccc}
\toprule
Window (quarters) & Mean $\hat{\rho}$ & Std.\ Dev. & Min & Max \\
\midrule
32 & 0.268 & 0.331 & $-$0.413 & 0.867 \\
40 & 0.290 & 0.344 & $-$0.354 & 0.859 \\
48 & 0.296 & 0.349 & $-$0.358 & 0.867 \\
\bottomrule
\end{tabular}
\begin{tablenotes}
\small
\item \textit{Code reference:} \texttt{code/empirical/appendix\_D/window\_sensitivity.py}
\end{tablenotes}
\end{threeparttable}
\end{table}

\subsection{Appendix E: MAB vs.\ Bayesian Learning}

Both microfoundations produce nearly identical aggregate dynamics.

\begin{table}[htbp]
\centering
\caption{Figure E1: MAB vs.\ Bayesian Comparison}
\label{tab:mab_bayesian}
\begin{threeparttable}
\begin{tabular}{lc}
\toprule
Metric & Value \\
\midrule
Correlation ($\theta_t$ vs.\ $\bar{w}_t$) & $>$0.87 \\
RMSE (inflation paths) & $<$0.25 pp \\
Asymmetry preserved & Yes (both frameworks) \\
\bottomrule
\end{tabular}
\begin{tablenotes}
\small
\item \textit{Code reference:} \texttt{code/simulations/appendix\_E/mab\_vs\_bayesian.py}, \texttt{code/plotting/appendix\_E/figure\_E1\_mab\_vs\_bayesian.py}
\item \textit{Output:} \texttt{output/figures/appendix\_E/figure\_E1\_mab\_vs\_bayesian.png}
\end{tablenotes}
\end{threeparttable}
\end{table}

%%%%%%%%%%%%%%%%%%%%%%%%%%%%%%%%%%%%%%%%%%%%%%%%%%%%%%%%%%%%%%%%%%%%%%%%%%%%%%%
\section{Computational Details}
\label{sec:computational}

\subsection{Algorithm Parameters}

The sophisticated FIRE solution algorithm uses the following configuration:

\begin{table}[htbp]
\centering
\caption{Numerical Algorithm Parameters}
\label{tab:algorithm_params}
\begin{tabular}{llc}
\toprule
Parameter & Description & Value \\
\midrule
$\alpha$ & Dampening factor & 0.3 \\
tol & Convergence tolerance & $10^{-8}$ \\
max\_iter & Maximum iterations & 200 \\
Typical convergence & Observed iterations & 40--100 \\
Runtime & $T = 120$ quarters & $<$1 second \\
\midrule
$\gamma$ & Softmax inverse temperature & 25,000 \\
Transition width & Softmax smoothing & $\pm$0.0002 \\
Error vs.\ hard indicator & Softmax approximation & $<$0.1\% \\
\bottomrule
\end{tabular}
\end{table}

%%%%%%%%%%%%%%%%%%%%%%%%%%%%%%%%%%%%%%%%%%%%%%%%%%%%%%%%%%%%%%%%%%%%%%%%%%%%%%%
\section{Summary of Key Numerical Results}
\label{sec:summary}

\begin{table}[htbp]
\centering
\caption{Summary: Key Numerical Claims in Manuscript}
\label{tab:summary}
\begin{threeparttable}
\begin{tabular}{lll}
\toprule
Claim & Value & Code Reference \\
\midrule
\multicolumn{3}{l}{\textit{Section 2: Model}} \\
Credibility recovery asymmetry & $\approx$15 years & \texttt{simulations/section\_2/shock\_regimes.py} \\
Persistence triggers de-anchoring & $\rho = 0.85$ & \texttt{tables/section\_2/table\_02\_transitory\_persistent.py} \\
\midrule
\multicolumn{3}{l}{\textit{Section 3: Applications}} \\
Oil shock peak difference & 2.7 pp & \texttt{simulations/section\_3/oil\_shocks.py} \\
Oil shock persistence ratio & $3.0\times$ & \texttt{simulations/section\_3/oil\_shocks.py} \\
Great Moderation correlation & 0.94 & \texttt{simulations/section\_3/great\_moderation.py} \\
Missing disinflation: NK vs.\ Adaptive & $-3.7\%$ vs.\ $-0.6\%$ & \texttt{simulations/section\_3/missing\_disinflation.py} \\
Post-pandemic: NK vs.\ Adaptive peak & 16.7\% vs.\ 7.8\% & \texttt{simulations/section\_3/post\_pandemic.py} \\
Transmission lag difference & 40\% & \texttt{simulations/section\_3/transmission\_lags.py} \\
Welfare loss ratio (low vs.\ high $\theta$) & $3.6\times$ & \texttt{simulations/section\_3/policy\_asymmetry.py} \\
\midrule
\multicolumn{3}{l}{\textit{Section 4: Extensions}} \\
Brazil welfare improvement & $>$90\% & \texttt{simulations/section\_4/brazil.py} \\
Kalman vs.\ rolling correlation & 0.86 & \texttt{empirical/appendix\_D/kalman\_comparison.py} \\
MAB vs.\ Bayesian correlation ($	heta_t$ vs.\ $ar{w}_t$) & $>$0.87 & \texttt{simulations/appendix\_E/mab\_vs\_bayesian.py} \\
\bottomrule
\end{tabular}
\begin{tablenotes}
\small
\item All code files located in \texttt{code/simulations/}, \texttt{code/plotting/}, \texttt{code/tables/}, or \texttt{code/empirical/}.
\end{tablenotes}
\end{threeparttable}
\end{table}

\end{document}
